% =====================================================================
%  Course syllabus                                       (English edition)
% =====================================================================
\documentclass[11pt,a4paper]{article}
\providecommand\ACNroot{../..}
\input{\ACNroot/templates/lecture-en}

\ACNset{number}{0}
\ACNset{title}{Course Syllabus}
\ACNset{shorttitle}{Syllabus}

\begin{document}
\begin{center}
  {\sffamily\bfseries\color{acnnavy}\LARGE Advanced Communication Networks\par}
  \vspace{2pt}
  {\sffamily\color{acnteal}\Large Course Syllabus\par}
  \vspace{2pt}
  {\fa{شبکه‌های ارتباطی پیشرفته -- طرح درس}\par}
\end{center}

\section{Course description}
This course follows mobile communication networks from the second to the sixth generation. It explains the radio access, core network, signalling,
mobility, security and quality-of-service mechanisms of GSM/GPRS, UMTS, LTE/EPC and 5G NR/5GC, and the software-driven technologies (SDN, NFV,
service-based architecture) that transformed them. Each lecture combines architecture, procedures and quantitative models with worked examples and
graded exercises. All material is available in English and Persian with identical structure, notation, figures and numbering.

\section{Learning outcomes}
On completing the course students will be able to:
\begin{objectives}
  \item \textbf{Describe} the architecture and functions of 2G, 3G, 4G and 5G networks and map functions across generations.
  \item \textbf{Trace} the main procedures (attach/registration, authentication, session setup, handover, paging) in sequence diagrams.
  \item \textbf{Apply} quantitative models -- traffic and queueing, link budgets, scheduling, capacity, latency and signalling load -- to dimension networks.
  \item \textbf{Explain} SDN, NFV and the service-based architecture and their role in 5G slicing and QoS.
  \item \textbf{Read} 3GPP and ITU specifications and \textbf{assess} emerging 6G technologies critically.
\end{objectives}

\section{Prerequisites}
Probability and basic queueing, signals and systems, digital communications (modulation, Shannon capacity) and computer networking (TCP/IP).

\section{Course structure}
The 26 lectures are grouped into five parts. Each lecture is a standalone document with fourteen sections: objectives, prerequisites,
introduction, fundamentals, technical detail, mathematical models, architecture, procedures, worked examples, technology comparison, key takeaways,
exercises at three levels, solutions and further reading.

\begin{table}[htbp]
\centering\small
\caption{Course modules.}\label{tab:modules}
\rowcolors{2}{acnmist}{white}
\begin{tabularx}{\linewidth}{@{}r Y l@{}}
\toprule
\rowcolor{white} \# & Lecture & Generation \\
\midrule
1 & GSM, GPRS and EDGE & 2G \\
2 & 2G Circuit Switching & 2G \\
3 & 2G/3G Core Networks & 3G \\
4 & Mobility and Handover & 4G \\
5 & 4G LTE Architecture & 4G \\
6 & E-UTRAN: the LTE Radio Access Network & 4G \\
7 & LTE Scheduling & 4G \\
8 & Random Access Channel (RACH) & 4G \\
9 & Evolved Packet Core (EPC) & 4G \\
10 & Diameter Protocol & 4G \\
11 & Home Subscriber Server (HSS) & 4G \\
12 & Mobility Management Entity (MME) & 4G \\
13 & Serving and PDN Gateways (S-GW and P-GW) & 4G \\
14 & 5G Networks: Introduction & 5G \\
15 & Software-Defined Networking (SDN) & 5G \\
16 & Network Functions Virtualisation (NFV) & 5G \\
17 & 3GPP Standardisation & 5G \\
18 & 5G Service-Based Architecture (SBA) & 5G \\
19 & 5G Core Network Architecture & 5G \\
20 & 5G System Procedures & 5G \\
21 & Network Slicing & 5G \\
22 & Quality of Service (QoS) in 5G & 5G \\
23 & The 5G User Plane & 5G \\
24 & The 5G Control Plane & 5G \\
25 & Network Evolution: a Grand Synthesis & 5G \\
26 & Future 6G Technologies & 6G \\
\bottomrule
\end{tabularx}
\end{table}

\begin{itemize}
  \item \textbf{Part I -- 2G/3G foundations} (Lectures 1--4).
  \item \textbf{Part II -- 4G access and core} (Lectures 5--13).
  \item \textbf{Part III -- Software-driven architectures} (Lectures 15--18).
  \item \textbf{Part IV -- 5G system and procedures} (Lectures 14, 19--22).
  \item \textbf{Part V -- Planes and outlook} (Lectures 23--26).
\end{itemize}

\section{Suggested schedule}
For a 13-week semester, two lectures per week in numerical order; each lecture needs about 6--8 hours of study (3 hours of reading and 3--5 hours of
exercises). Weeks 1--2 cover Part I, weeks 3--7 Part II and Lecture 14, weeks 8--9 Part III, weeks 10--11 the rest of Part IV, and weeks 12--13 Part V.

\section{Suggested assessment}
The following weighting is a suggestion that instructors may adapt: weekly exercise sets 30\,\%, a mid-term examination on Parts I--II 25\,\%, a
project (for example, dimensioning or simulating a part of a 5G network) 15\,\%, and a final examination on the whole course 30\,\%.

\section{Materials}
The lecture notes, figures and bibliography of this repository; the 3GPP and ITU documents cited in each lecture; and the textbooks listed under
further reading, in particular Dahlman et al.\ (NR), Sesia et al.\ (LTE) and Rommer et al.\ (5G core).

\section{Language}
Both editions are equivalent; students may study in either language. Technical abbreviations are kept in their standard Latin form in both.
\end{document}
